% !TEX program = xelatex
\documentclass[14pt]{extarticle}

% ============================================================
% GEOMETRY
% ============================================================
\usepackage[
  a4paper,
  left=2.5cm,
  right=1.5cm,
  top=2cm,
  bottom=2cm
]{geometry}

% ============================================================
% МОВА ТА ШРИФТИ
% ============================================================
\usepackage[ukrainian]{babel}
\usepackage{fontspec}
\setmainfont{Times New Roman}

% ============================================================
% ОСНОВНИЙ ТЕКСТ
% ============================================================
\usepackage{setspace}
\onehalfspacing
\setlength{\parindent}{1.27cm}
\setlength{\parskip}{0pt}
\usepackage{microtype}

% ============================================================
% ЗАГОЛОВКИ
% ============================================================
\usepackage{titlesec}

\titleformat{\section}
  {\bfseries\fontsize{16pt}{24pt}\selectfont\centering}
  {\thesection}{0.5em}{}
\titlespacing*{\section}{0pt}{20pt}{6pt}

\titleformat{\subsection}
  {\bfseries\fontsize{14pt}{21pt}\selectfont}
  {\thesubsection}{0.5em}{}
\titlespacing*{\subsection}{1.27cm}{6pt}{6pt}

\titleformat{\subsubsection}
  {\bfseries\itshape\fontsize{14pt}{21pt}\selectfont}
  {\thesubsubsection}{0.5em}{}
\titlespacing*{\subsubsection}{1.27cm}{6pt}{0pt}

% ============================================================
% СПИСКИ
% ============================================================
\usepackage{enumitem}

\setlist[itemize,1]{
  label=$\bullet$,
  leftmargin=1.27cm,
  labelwidth=0.5cm,
  labelsep=0.3cm,
  itemindent=0pt,
  itemsep=3pt,
  topsep=6pt,
  parsep=0pt,
  partopsep=0pt
}

\setlist[enumerate,1]{
  label=\arabic*.,
  leftmargin=2.54cm,
  labelwidth=0.63cm,
  labelsep=0.54cm,
  itemindent=0pt,
  itemsep=0pt,
  topsep=0pt,
  parsep=0pt,
  partopsep=0pt
}

% ============================================================
% ТАБЛИЦІ
% ============================================================
\usepackage{array}
\usepackage{longtable}
\usepackage{booktabs}

% ============================================================
% ГІПЕРПОСИЛАННЯ
% ============================================================
\usepackage[hidelinks,breaklinks=true]{hyperref}

\begin{document}

\sloppy
\emergencystretch=3em

% ==================== ТИТУЛЬНА ЧАСТИНА ====================
\begin{center}
\textbf{Лабораторна робота 1}

\medskip

\textbf{\MakeUppercase{Визначення сфери застосування (комерційності) проєкту на основі 7 базових питань. Визначення цілей проєкту за системою SMART.}}
\end{center}

\medskip

\textbf{Мета роботи:} ознайомитись з принципами визначення характеристик проєкту на основі 7 базових питань. Навчитись визначити цілей проєкту за системою SMART.

\textbf{Проєкт:} <<Вдосконалення токенізатора української мови для великих мовних моделей>>.

\textbf{Виконав:} студент групи ШІ-42 Мирош Андрій.

% ==================== 7 БАЗОВИХ ПИТАНЬ ====================
\section*{1. Визначення сфери застосування на основі 7 базових питань}
\addcontentsline{toc}{section}{1. Визначення сфери застосування на основі 7 базових питань}

Метод <<Критичних запитань>> передбачає послідовну відповідь на 7 ключових запитань щодо проєкту, які дозволяють визначити його доцільність, межі та зацікавлені сторони.

\subsection*{1.1. Хто?}

Зацікавленими сторонами проєкту є:
\begin{itemize}
  \item \textbf{Розробники LLM та NLP-систем} --- інженери, які адаптують великі мовні моделі до української мови і потребують ефективного токенізатора, що зменшує довжину токен-послідовностей та покращує використання контекстного вікна;
  \item \textbf{Дослідники NLP} --- науковці, які порівнюють токенізатори різних LLM на українських корпусах і потребують стандартизованого інструментарію та метрик оцінювання;
  \item \textbf{Розробники україномовних сервісів} --- команди, які створюють чат-ботів, пошукові системи та інші текстові продукти для української мови і хочуть обрати оптимальний токенізатор;
  \item \textbf{Науковий керівник та кафедра} --- академічна зацікавленість у результатах дослідження ефективності токенізації;
  \item \textbf{Виконавець} --- студент групи ШІ-42 Мирош Андрій, який реалізує проєкт у рамках бакалаврської кваліфікаційної роботи.
\end{itemize}

\subsection*{1.2. Що?}

Проєкт передбачає вдосконалення (розробку та налаштування) спеціалізованого BPE-токенізатора української мови та побудову відтворюваного конвеєру порівняння його якості з токенізаторами сучасних LLM. Конкретно створюються:
\begin{itemize}
  \item локальний BPE-токенізатор зі словником 32\,000 одиниць, навчений на українських текстах;
  \item уніфікована процедура попередньої обробки тексту (Unicode NFC-нормалізація, уніфікація апострофів);
  \item набір метрик оцінювання: token fertility (середнє та p95), частка UNK/OOV, частка byte fallback, round-trip fidelity;
  \item MVP-система з модульною архітектурою (завантаження корпусу, токенізація, обчислення метрик);
  \item CLI-інтерфейс для інтерактивного тестування токенізатора та запуску обчислень на корпусі.
\end{itemize}

\subsection*{1.3. Де?}

Проєкт реалізується на кафедрі систем штучного інтелекту Національного університету <<Львівська політехніка>>. Розробка ведеться локально з використанням Python-екосистеми (Hugging Face Transformers, SentencePiece). Обчислення не потребують GPU --- всі експерименти працюють лише з токенізацією (<<текст $\to$ токени>> та <<токени $\to$ текст>>), без запуску повноцінних мовних моделей. Код та результати розміщені у відкритому репозиторії на GitHub.

\subsection*{1.4. Коли?}

Проєкт обмежений термінами виконання бакалаврської кваліфікаційної роботи (один навчальний семестр). Етапи виконання:
\begin{enumerate}
  \item Аналіз літературних джерел та формування вимог до токенізації українських текстів.
  \item Аналіз та вибір корпусу для порівняння (Kobza Corpus).
  \item Реалізація/адаптація локального BPE-токенізатора та уніфікованої попередньої обробки.
  \item Визначення та реалізація метрик (token fertility, UNK/OOV, byte fallback, round-trip fidelity).
  \item Експериментальне порівняння з токенізаторами 5 поширених LLM на 100\,000 зразках корпусу.
  \item Аналіз результатів, оформлення та підготовка до захисту.
\end{enumerate}

\subsection*{1.5. Чому?}

Токенізація є одним із ключових вузлів у пайплайні сучасних LLM: вона визначає спосіб перетворення тексту на дискретні індекси словника та безпосередньо впливає на довжину послідовностей у токенах, використання контекстного вікна, швидкодію та обчислювальні витрати. Для української мови, яка характеризується розвиненою словозміною, продуктивним словотворенням та специфічними орфографічними особливостями (наприклад, варіативність апострофа), стандартні <<англоцентричні>> токенізатори часто дають неефективне представлення --- довші послідовності токенів на слово.

Незважаючи на зростання кількості українських NLP-ресурсів, стандартизованого інструментарію для порівняння токенізаторів саме на українських текстах не існує. Тому створення спеціалізованого токенізатора та відтворюваного конвеєру оцінювання є актуальним завданням.

\subsection*{1.6. Як?}

Підхід складається з таких кроків:
\begin{enumerate}
  \item \textbf{Підготовка даних} --- завантаження корпусу Kobza з Hugging Face у потоковому режимі (streaming), застосування єдиної нормалізації (Unicode NFC + уніфікація апострофів), сегментація слів за пробілами.
  \item \textbf{Токенізація} --- уніфікований інтерфейс для локального BPE-токенізатора та токенізаторів LLM (через Hugging Face AutoTokenizer); контроль відсутності службових токенів (BOS/EOS) для чистоти метрик.
  \item \textbf{Обчислення метрик} --- автоматичний прогін корпусу кожним токенізатором, обчислення token fertility (середнє, p95), частки UNK, частки byte fallback та round-trip fidelity.
  \item \textbf{Порівняння} --- 6 токенізаторів (Local BPE, aya-expanse-8b, gemma-3-12b-it, Llama-3.1-8B, Phi-4-mini, Qwen2.5-7B) на 100\,000 зразках; збереження результатів у структурованому форматі для візуалізації.
\end{enumerate}

\subsection*{1.7. Скільки?}

Проєкт має один основний результат --- MVP-система порівняння токенізаторів. Кількісні характеристики:
\begin{itemize}
  \item 1 локальний BPE-токенізатор зі словником 32\,000 одиниць;
  \item порівняння з 5 токенізаторами поширених LLM;
  \item оцінювання на 100\,000 зразках корпусу Kobza;
  \item 4 метрики оцінювання (token fertility mean/p95, UNK rate, byte fallback rate, round-trip fidelity);
  \item 1 CLI-інструмент для тестування;
  \item відкритий репозиторій із кодом та результатами.
\end{itemize}

\subsection*{1.8. Як ми знаємо, що досягли цілей?}

Успішність проєкту визначається за такими критеріями:
\begin{itemize}
  \item локальний BPE-токенізатор навчений і демонструє нижчу середню token fertility порівняно з токенізаторами LLM (ціль: $<$2 токени/слово);
  \item метрики обчислені для всіх порівнюваних токенізаторів на одному корпусі з однаковою попередньою обробкою;
  \item результати відтворювані: запуск конвеєру на тому ж корпусі дає ідентичні значення метрик;
  \item CLI-інтерфейс функціонує та дозволяє токенізувати довільний текст і запускати обчислення метрик;
  \item виявлені сильні та слабкі сторони кожного токенізатора задокументовані.
\end{itemize}

\subsection*{1.9. Що виходить за рамки?}

За межами проєкту залишаються:
\begin{itemize}
  \item навчання або запуск повноцінних мовних моделей --- оцінюється лише токенізація (<<текст $\to$ токени>> та <<токени $\to$ текст>>);
  \item оцінювання впливу токенізатора на downstream-задачі (генерація, класифікація тощо);
  \item розробка нових алгоритмів субсловної сегментації --- використовується стандартний BPE;
  \item оптимізація продуктивності/швидкості токенізатора;
  \item покриття інших мов, окрім української;
  \item комерційне розгортання або інтеграція у виробничі системи.
\end{itemize}

% ==================== ТРИКУТНИК ОБМЕЖЕНЬ ====================
\subsection*{1.10. Трикутник конкурентних обмежень}

Після визначення сфери застосування можна заповнити трикутник конкурентних обмежень (час, ресурси, якість):

\begin{itemize}
  \item \textbf{Час} --- проєкт обмежений термінами виконання бакалаврської кваліфікаційної роботи (один навчальний семестр). Ключовий ризик --- затримка на етапі інтеграції різних токенізаторів через відмінності в API та форматах.

  \item \textbf{Ресурси} --- обчислювальні витрати мінімальні: усі експерименти працюють на CPU, оскільки виконується лише токенізація, а не інференс моделей. Основний ресурс --- час розробника. Дані доступні безкоштовно (Kobza Corpus на Hugging Face).

  \item \textbf{Якість} --- цільова якість визначається відтворюваністю конвеєру та інформативністю метрик. Ключовий показник --- token fertility: ціль для локального токенізатора $<$2 токени/слово. Також важлива коректність round-trip fidelity (ціль: 100\%).
\end{itemize}

\textbf{Пріоритетність:} у контексті бакалаврської роботи пріоритет надається якості порівняння та відтворюваності результатів, потім --- дотриманню термінів, і останнім --- обсягу покриття (кількість токенізаторів та корпусів).

% ==================== SMART ====================
\section*{2. Визначення цілей проєкту за системою SMART}
\addcontentsline{toc}{section}{2. Визначення цілей проєкту за системою SMART}

Метод SMART дозволяє оцінити мету проєкту за п'ятьма критеріями: конкретність (Specific), вимірюваність (Measurable), наявність виконавця (Assignable), реалістичність (Realistic), обмеженість у часі (Time-related).

\subsection*{2.1. S --- Specific (Конкретність)}

Мета проєкту сформульована конкретно: \textit{вдосконалити спеціалізований BPE-токенізатор української мови та побудувати відтворюваний конвеєр порівняння його якості з токенізаторами сучасних LLM за набором стандартизованих метрик на українському корпусі}.

Конкретність забезпечується чітким визначенням:
\begin{itemize}
  \item що саме створюється --- локальний BPE-токенізатор (32K словник) та MVP-система оцінювання;
  \item з чим порівнюється --- з токенізаторами 5 поширених LLM (aya-expanse-8b, gemma-3-12b-it, Llama-3.1-8B, Phi-4-mini, Qwen2.5-7B);
  \item на яких даних --- корпус Kobza (100\,000 зразків);
  \item за якими критеріями --- token fertility, UNK/byte fallback, round-trip fidelity;
  \item у якому форматі --- CLI-інструмент + структуровані результати для візуалізації.
\end{itemize}

\subsection*{2.2. M --- Measurable (Вимірюваність)}

Результати проєкту є вимірюваними за допомогою чітко визначених метрик:
\begin{itemize}
  \item \textbf{Token fertility (середнє)} --- середня кількість токенів на слово; ціль для локального токенізатора $<$2 (досягнутий результат: 1.803);
  \item \textbf{Token fertility (p95)} --- 95-й перцентиль фертильності; ціль $\leq$5 (досягнутий результат: 4);
  \item \textbf{Частка UNK} --- відсоток невідомих токенів; ціль 0\%;
  \item \textbf{Частка byte fallback} --- відсоток <<сирих>> байтових токенів;
  \item \textbf{Round-trip fidelity} --- частка рядків, що точно відтворюються після encode $\to$ decode; ціль 100\%.
\end{itemize}

Усі метрики обчислюються автоматично на одному корпусі з однаковою попередньою обробкою, що забезпечує об'єктивність порівняння.

\subsection*{2.3. A --- Assignable (Наявність виконавця)}

Проєкт виконується студентом групи ШІ-42 Мирошем Андрієм у рамках бакалаврської кваліфікаційної роботи за спеціальністю 122 <<Комп'ютерні науки>> під керівництвом Яковини В.\,С. Виконавець має необхідні компетенції у галузі обробки природної мови та програмування мовою Python.

Розподіл відповідальності:
\begin{itemize}
  \item \textbf{Виконавець} --- реалізація токенізатора, побудова конвеєру оцінювання, проведення експериментів, аналіз результатів, створення CLI, оформлення звіту;
  \item \textbf{Науковий керівник} --- методичне спрямування, контроль якості дослідження, рецензування.
\end{itemize}

\subsection*{2.4. R --- Realistic (Реалістичність)}

Реалістичність проєкту обґрунтовується такими факторами:
\begin{itemize}
  \item \textbf{Наявність даних} --- корпус Kobza є загальнодоступним на Hugging Face, охоплює різні домени та стилі мовлення, репрезентує реальну українську мову;
  \item \textbf{Перевірені методи} --- BPE є стандартним підходом до субсловної токенізації, підтримується SentencePiece; Hugging Face AutoTokenizer забезпечує єдиний інтерфейс до токенізаторів LLM;
  \item \textbf{Мінімальні обчислювальні вимоги} --- усі експерименти працюють на CPU, не потребують GPU, що робить проєкт доступним для виконання на будь-якому комп'ютері;
  \item \textbf{Наявність попередніх досліджень} --- існують роботи з оцінювання ефективності токенізації для української на токенізаторах foundational-моделей, що підтверджують актуальність підходу та надають орієнтири для порівняння;
  \item \textbf{Відкритість} --- весь код розміщено у відкритому репозиторії, що забезпечує відтворюваність.
\end{itemize}

\subsection*{2.5. T --- Time-related (Обмеженість у часі)}

Проєкт обмежений термінами виконання бакалаврської кваліфікаційної роботи. Основні етапи:
\begin{enumerate}
  \item Аналіз літератури та формування вимог до токенізації.
  \item Вибір корпусу та реалізація попередньої обробки.
  \item Реалізація локального BPE-токенізатора та уніфікованого інтерфейсу.
  \item Реалізація метрик та конвеєру порівняння.
  \item Проведення експериментів на корпусі Kobza (100\,000 зразків).
  \item Аналіз результатів, оформлення та підготовка до захисту.
\end{enumerate}

Кожен етап має чітку послідовність та залежності. Наявність кінцевої дати (захист бакалаврської роботи) забезпечує фокусування зусиль.

% ==================== ВИСНОВКИ ====================
\section*{Висновки}
\addcontentsline{toc}{section}{Висновки}

У ході виконання лабораторної роботи було обґрунтовано доцільність створення спеціалізованого BPE-токенізатора української мови та конвеєру порівняння за двома методами.

\textbf{Метод 7 базових питань} дозволив визначити: зацікавлені сторони (розробники LLM, дослідники NLP, розробники україномовних сервісів), суть проєкту (BPE-токенізатор + MVP-система порівняння з CLI), місце реалізації (НУ <<Львівська політехніка>>, CPU-based експерименти), часові рамки (навчальний семестр), обґрунтування актуальності (неефективність <<англоцентричних>> токенізаторів для української), спосіб реалізації (BPE + уніфікований інтерфейс + набір метрик), кількісні показники та критерії успіху. Також визначено межі проєкту (лише токенізація, без downstream-задач) та заповнено трикутник обмежень із пріоритетом на якість та відтворюваність.

\textbf{Метод SMART} підтвердив, що мета проєкту є конкретною (BPE-токенізатор 32K + порівняння з 5 LLM-токенізаторами на Kobza), вимірюваною (token fertility $<$2, p95 $\leq$5, UNK 0\%, round-trip 100\%), забезпеченою виконавцем (студент під керівництвом наукового керівника), реалістичною (наявність даних, CPU-based, перевірені методи) та обмеженою у часі (семестр виконання бакалаврської роботи).

Обидва методи свідчать про доцільність реалізації проєкту.

\end{document}
